\section{Benchmark Details}
\label{app:bench-details}

\subsection{Six-axis schema}
\begin{center}\small
\begin{tabularx}{\linewidth}{@{}llX@{}}
\toprule
ID                       & Capability tested                & Primary metrics \\
\midrule
\texttt{b1\_forecast}    & Temporal graph prediction        & PageRank/HHI/Gini MAE, trend consistency, Spearman; per $h\!\in\!\{1,4,8,12\}$ \\
\texttt{b2\_warning}     & Streaming anomaly detection      & Precision, recall, F1, lead time, alert stability; per event $\times$ budget \\
\texttt{b3\_calibration} & Predictive uncertainty           & PI coverage at $0.90$ target, PI width, ECE, CRPS \\
\texttt{b4\_stress}      & What-if scenario fidelity        & Loss MAE, distressed-count MAE, propagation depth MAE, target overlap@$k$ \\
\texttt{b5\_decision}    & Supervisory ticket quality       & Precision, recall@$k$, F1, target stability, judge score, FIR \\
\texttt{b6\_robustness}  & Data quality sensitivity         & Per-benchmark metrics under 5 degradation regimes; relative degradation \\
\bottomrule
\end{tabularx}
\end{center}

The six axes are chosen to be (i)~\emph{non-overlapping}---a system
can excel at \texttt{b1} while failing \texttt{b5};
(ii)~\emph{decomposable along the four-layer framework}, so that
ablations on the predictor layer (\texttt{b1}, \texttt{b3},
\texttt{b6}), monitor layer (\texttt{b2}), scenario layer
(\texttt{b4}), and decision layer (\texttt{b5}) can be reported
independently; and (iii)~\emph{individually publishable}, so that
downstream work may extend a single axis without re-running the
whole suite.

\subsection{Dataset split}
\begin{center}\small
\begin{tabularx}{\linewidth}{@{}llrX@{}}
\toprule
Split      & Period             & Weeks       & Usage \\
\midrule
Train      & 2020-03--2024-06   & $\sim$222   & Forecaster training \\
Validation & 2024-07--2024-12   & $\sim$26    & Conformal calibration, threshold tuning \\
Test       & 2025-01--2025-08   & $\sim$33    & Frozen leaderboard \\
\bottomrule
\end{tabularx}
\end{center}

\subsection{Ground-truth equation}
With $w_t(v)\!=\!\sum_{e\ni v}w_t(e)$ the total weight incident on
$v$, the regulator-aligned stressed set at horizon $h$ is
\begin{equation}
\Delta_t^h(v)\!=\!w_t(v)\!-\!w_{t+h}(v),\quad
\mathcal{S}_t^h \!=\! \mathrm{top}_{\pi}\bigl\{ v : w_t(v)\!>\!0,\;
                                              \Delta_t^h(v)\!>\!0 \bigr\},
\label{eq:stressed-app}
\end{equation}
with $\pi\!=\!0.05$ and ground-truth horizon $h\!=\!4$~weeks. The
final four test weeks lack the $4$-week-ahead snapshot $w_{t+4}$
and are not scored, leaving $n\!=\!29$ evaluated weeks. This is the
single ground-truth definition used by \texttt{b2}, \texttt{b5}, and
the LLM-decision evaluation pipeline.

\subsection{Reporting protocol}
A submission produces one JSON file per
\{benchmark,~method\} pair following the harness schema, plus a
single \texttt{results.json} aggregating the suite. The full
leaderboard requires \texttt{b1}, \texttt{b3}, \texttt{b5},
\texttt{b6}; \texttt{b2} and \texttt{b4} are required only for
systems that implement a monitor or scenario layer respectively. The
harness fixes random seeds, snapshots library versions, and emits
SHA-256 hashes of all checkpoint files.
